\section{PCB}
\label{sec:pcb}

我的进程/线程信息分为两块：TCB (\textit{thread control block}) 和 PCB (\textit{process control block})。调度器中存储的都是 TCB。

TCB 中的内容是每个线程单独的信息，例如线程当前是否阻塞，\cref{sec:interrupt}中介绍的 ksp 与 usp，以及\cref{subsec:syscallctx}中提到的 system-call context。它还含有一个指向线程共有的 PCB 的指针。

至于 PCB，粗略来看有这些东西：
\begin{itemize}
\item pid;
\item uid, suid, euid（以及对应的 gid, sgid, egid）;
\item 进程开启时的 PC;
\item 父进程与子进程；
\item VMA (\textit{virtual memory area});
\item 开启的文件，以及 VFS context；
\item 等待处理的信号，以及 signal handler；
\item 进程的返回值。
\end{itemize}

值得注意的是，在启动时，我会生成一个 pid = -1 的假 PCB，作为调度器的“当前进程”，并在启动完成后移除。这可以简化大量启动时的特殊判断。

\subsection{初始化}
\label{subsec:pcbinit}

进程分为两种：kernel- 和 user-space。

我们先考虑 user-space 进程。ELF loader 将会根据 ELF 的内容，填充 PCB 的 VMA，但暂时不将其 map 到页表中。在程序真正需要访问这些地址的时候，触发 page-fault，才会真正地分配物理内存。实际上，fork() 还支持 copy-on-write，具体的实现方法是禁用这个进程和子进程的所有写权限，在需要写的时候重新 map。有 VMA 在，不需要做多少处理就可以实现 COW 了。对于 VMA 的更详细说明，请见\cref{subsec:vma}。

说到页表，user-space 的页表分为两部分：下半部分直接 shallow-copy 内核的页表，而上半部分则是独属于这个进程的，最开始则是空的。在\cref{sec:boot}中提到过，Sv39 的页表被标准强行分为两部分，我的 OS 处于高位部分，而 user-space 则使用低位部分。这使得所有的进程都可以共用半个页表，而且 syscall 的时候不需要再重新设置 satp 寄存器（页表的根部）了。

在 ELF loader 完成之后，我们还会给它分配一个堆和一个栈。堆的初始大小只有一页，栈则是 8 MB。

对于 kernel-space 的进程，初始化就简单很多，直接往对应的函数一指就可以了。堆就靠整个 kernel 共用的那个分配器，不过栈依然是单独的。

最后，打开 stdin/stdout/stderr 三个 fd，就初始化完成了。注意如果没有找到 /dev/console，这个 OS 会直接 panic。

\subsection{System-call context}
\label{subsec:syscallctx}

一部分 system-call 可能会 \textit{block}，也就是让当前进程暂停执行。我们需要保存当前的执行进度，然后立即切换到下一个进程。PCB 中的 system-call context 字段就是用来干这个的。

正如\cref{sec:interrupt}中介绍的中断处理一样，保存“进度”实际上只需要保存合适的寄存器。实际上，由于暂停执行实际上是由 \cpp{suspend()} 函数\footnote{这其实是个宏，定义为\cpp{context_save(&ctx, &ctx_valid)}。它本质上还是一个函数，所以后面的内容依旧适用。}实现，编译器不会期待我们保存 caller-saved registers。所以，我们需要保存的就只有 s0-s11，以及从暂停中恢复时所执行的“下一条”指令。观察到这就是 \cpp{suspend} 的返回地址，所以我们保存 ra 就可以了。在保存完毕之后，我们调用 \cpp{scheduler.dispatch()}，直接切换到下一个进程。

保险起见，我依然保留了 sstatus 和 sepc。虽然它们在 ksp 那里被保存得很好，但我们可能会在 syscall 的时候直接读取这些寄存器的值，而不是去 trapframe 里找，所以最好还是保存一下。

一个陷阱是，我们不能标注 \cpp{[[noreturn]] void suspend()}。这是因为 suspend 虽然从实现上来说不会返回，但它其实最终还是会“跳回来”的，所以从这个角度，它和普通函数无异。与之相对的是 \cpp{context_restore} 是可以标注为 noreturn 的。

\subsection{VMA}
\label{subsec:vma}

当分配一段内存的时候，我们显然需要保存这段内存的开始和结束地址。不过，这并不是全部。

在加载 ELF 文件的时候，我们会遇到一些 PT\_LOAD 段，指明它们需要被加载进内存。我们并不想把它一次性全加载进去，而是按需加载。因此，在这样分配 VMA 的时候，还必须记录它的 ``backup file'' 以及偏移。如果文件不够长，那么我选择在剩余位置补零：这可以很好地实现 .bss 段。

在处理 brk 系统调用的时候，我们希望扩大堆的大小，那么我们必须得从一串 VMA 中找到哪一块是堆，并修改它的结束地址。当然，另一种解决方案是直接再增加一块 VMA，但创建这块新的 VMA 时，依然需要知道它的起始地址——也就是堆的结束地址。所以，我给每个 VMA 都增加了一个标记，标明它是来自 ELF 的 PT\_LOAD 段，还是初始化时分配的堆或者栈。
